% ===============================================
% ============== Document setup =================
% ===============================================

% Specify your preferred document language, your type of
% thesis, and the professorship under which your thesis 
% takes place. To that end, use the following class options
% of our "vcethesis" document class:

% Language Options
%   english
%   german
%
% Thesis Type
%   bsc
%   msc
%
% Professorship / Chair
%
%   egger
%   guenther
%   stamminger
%   weyrich
%
% Example usages:
%
%   \documentclass[english, egger, bsc]{vcethesis}
%   \documentclass[german, stamminger, msc]{vcethesis}
%   \documentclass[english, weyrich, msc]{vcethesis}
%     ...
%
% Beyond that, unless you are directly supervised by 
% the professor, you will have to enter your direct
% supervisor. See below for further instructions.
%
% Now edit the following line as needed:
\documentclass[english, egger, bsc]{vcethesis}

% ===============================================
% ================= Packages ====================
% ===============================================
% Feel free to customize used packages. 
% See vcethesis.cls in the packages section, to see 
% the full list of included packages

\usepackage[textsize=small,textwidth=0.85\ML]{todonotes}

% ===============================================
% ================== Macros =====================
% ===============================================

% here are some macro definitions that you may find useful:
\input macros

% ===============================================
% ============== Configurations =================
% ===============================================
% Title of your thesis
\title{Counting Fruits with State-of-the-Art Multi-View-Stereo Approaches}        

% Author name
\author{Alexander Link}

% Birth day
\birthday{29.08.2004}

% Birth place
\birthplace{Bamberg}

% If your supervisor is not a professor, uncomment the following line to list the supervisor.
%\supervisor{<Supervisor>}
% for multiple supervisors, use \supervisors{<Supervisor1> and <Supervisor2>}
% for a german thesis, use \supervisorM{<Supervisor>}, \supervisorF{<Supervisor>} and \supervisorD{<Supervisor>} for m/f/d supervisor.

% Begin of thesis
\bdate{01.05.2026}

% End of thesis
\edate{01.10.2026}

% Additional logos: if your work is jointly supervised across
% chairs/professorships, uncomment the following line and add
% additional logos as needed:
%
%\additionallogos{\logoegger\logoweyrich\logostamminger}
%\additionallogos{\logoweyrich\logostamminger}

% Declaration of usage of AI tools:
% Usage of AI tools has to be explicitly declared in the thesis in the declaration section.
% See https://www.studon.fau.de/studon/go/copa/6188636 for more details.
% Please take the following formulation and adapt it to reflect your usage:
%
\aideclaration{During the preparation of this work, AI technologies were used to assist in the writing process. Specifically, <AI Tool 1 (Company, City, State, Country)> was used in order to check for grammar and style consistency, and <AI Tool 2 (Company, City, State, Country)> was used in order to assist with rephrasing and improving readability. After using these tools, the manuscript was carefully reviewed and the content was edited as needed. No tools or services were used for content generation. The content of this thesis was entirely created by myself, and I ensured that all material presented reflects my original work and research. I take full responsibility for the content of the thesis.}
%
% If you are writing your theses in german, use the following as a template.
%
% \aideclaration{Während der Vorbereitung dieser Arbeit wurden KI-Technologien eingesetzt, um den Schreibprozess zu unterstützen. Insbesondere wurde <AI Tool 1 (Company, City, State, Country)> verwendet, um die Konsistenz von Grammatik und Stil zu überprüfen, und <AI Tool 2 (Company, City, State, Country)> wurde eingesetzt, um bei der Umformulierung und Verbesserung der Lesbarkeit zu helfen. Nach dem Einsatz dieser Tools wurde das Manuskript sorgfältig überprüft und der Inhalt nach Bedarf überarbeitet. Für die Erstellung des Inhalts wurden keine Tools oder Dienste verwendet. Der Inhalt dieser Arbeit wurde vollständig von mir selbst erstellt, und ich habe sichergestellt, dass das gesamte dargestellte Material meine eigene Arbeit und Forschung widerspiegelt. Ich übernehme die volle Verantwortung für den Inhalt der Arbeit.}

% ===============================================
% ================= Document ====================
% ===============================================
\begin{document}

% -----------------------------------------
% ----------- Front Matter ----------------
% -----------------------------------------
\prepages
\maketitle
\cleardoublepage
% -----------------------------------------
% ------------- Abstract ------------------
% -----------------------------------------
% German: Abstrakt
\chapter*{Abstract}
%A brief and concise summary of the thesis containing the problem statement and how it was solved in about five sentences. After having read the abstract, a reader must know what the thesis is about.
------------
% ----------- TOC -------------
% -----------------------------------------
\tableofcontents
\cleardoublepage
% -----------------------------------------
% ----------- Main Document ---------------
% -----------------------------------------
\mainbody

% German: Einleitung
\chapter{Introduction}
%The introduction needs to answer the following questions:
%\begin{itemize}
%    \item Which area of visual computing are we in? (scientific visualization, appearance modeling, Monte Carlo rendering, etc.)
%    \item Which problem are we going to solve?
%    \item Why is this a challenging problem? What makes it difficult?
%    \item Who will be the user of this method? Who benefits from it?
%    \item How are you going to solve the problem?
%    \item How will you evaluate the method?
%    \item What are your most important findings?
%\end{itemize}
Feeding a growing world population is a major global challenge, and projections estimate that total food demand will rise substantially over the coming decades \cite{vandijk2021foodsecurity}. Precision agriculture responds to this by using advanced technology and data analysis to maximize crop yields while reducing waste and environmental impact \cite{karunathilake2023precision}. Within precision agriculture, counting the fruit on a tree is an important building block: it provides the yield estimates that growers rely on to plan the harvest and allocate labor, storage, and other resouces \cite{miranda2023}. Because performing the count manually across many trees is slow and labor-intensive, automating it from images is a valuable capability.

Among the ways to automate fruit counting, reconstructing the tree in 3D from a set of RGB images is particularly attractive, for two reasons. First, lifting the count into a shared three-dimensional representation lets each fruit be counted exactly once, instead of being counted repeatedly across the many overlapping images in which it appears. Second, the only hardware this requires is an ordinary color camera, of the kind build into any consumer device. This avoids the specialized and more expensive depth-sensing equipment, such as LiDAR scanners or calibrated stereo rigs, on which other 3D-caprure methods depend, and makes the approach inexpensive and easy to deploy in the field.
\todo[inline]{citation for cost, other 3D capture methods}

Recent work such as FruitNeRF \cite{meyer2024fruitnerf} shows that image-based 3D fruit counting is feasibly by training a Neural Radiance Field \cite{mildenhall2020nerf} per scene and counting the fruit in the resulting volumetric representation. This approach

\section{Challenges}

Counting fruits in this scenario is challenging for several reasons. Fruits are often obscured by foliage and by one another, the appearance of a scene varies with changing lighting, and the geometric complexity of a tree canopy makes a consistent reconstruction difficult. To reduce occlusion and the risk of missing apples, a scene should ideally be captured from all sides. However, the number of views is limited by hardware constraints such as available memory, so the cameras are sparsely distributed across a full 360-degree range, and widely spaced views often capture completely different regions of the tree. This makes it difficult to establish correspondences between views.

These conditions result in two opposing types of error. Individual apples observed by only a few cameras yield unstable correspondences and are easily split into multiple detections, whereas dense clusters of apples seen by many cameras may merge into a single detection. A single set of static parameters can therefore rarely resolve both cases optimally at the same time. Balancing these two failure modes is a central difficulty that our method must address.

\section{Our Approach}

To count the fruit in a scene, our method maps semantic information, in the form of 2D instance masks obtained with SAM3 \cite{carion2026sam3segmentconcepts}, across views by means of the 3D representation predicted by VGGT. The association of masks across views follows the pointmap-based principle introduced by CCGS \cite{hu2026ccgs}: pixel correspondence between views are established through the predicted pointmaps, and instance masks are matched across views by solvign an assignment problem that allows partial matches. We adopt this association step, transfer it from the image-pair setting of its origin to the multi-view reconstruction produced by VGGt, and build a new counting pipeline on top of it. In this pipeline, the pairwise mask matches are aggregated into multi-view instances through a graph formulation, and the fruit are counted by combining these associated instances with a subsequent clustering step.

\section{Contributions}
This thesis makes the following contributions. We investigate whetehr the feed-forward multi-view stereo model VGGT can replace per-scene optimization for image-based fruit counting, and we assess whether its predicted geometry is accurate enough for this task. Building on the pointmap-based mask association introduced by CCGS \cite{hu2026ccgs}, we develop a fruit-counting pipeline that extends this association with a graph-based aggregation of instances across many 360-degree-views, sphere fitting as an apple model, and a final clustering step to obtain the count. We further propose an iterative refinement of this pipeline that detects and resolves over-merged instances, thereby addressing the opposing split-and-merge errors that arise under sparse, widely spaced views. Finally, we evaluate both the non-iterative and the iterative variant on \todo{real-world data?}

% German: Verwandte Arbeiten
\chapter{Related Work}
\section{Fruit Counting and Yield Estimation}
Methods for counting fruit differ mainly in where the counting takes place: in the individual image, across a sequence of images, or in a three-dimensional representation of the scene. These stages can be read as increasingly robust answers to the core difficulty of counting across views, namely avoiding that the same fruit is counted more than once.

The earliest and simplest methods operate purely in the image plane, separating the task into detecting fruit in each image and counting the detections \cite{haeni2019}. Detection has moved from handcrafted color and shape features to deep neural networks \cite{miranda2023}, and modern detectors such as YOLO reach high accuracy on individual images \cite{ballenaruiz2025}. Because these methods count on a per-frame basis, however, they have no notion of which detections in different images correspond to the same fruit, so a fruit visible in several frames is counted several times. 

To resolve this, a second line of work tracks fruit across consecutive frames. Liu et al. \cite{liu2018} proposed the first complete pipeline of this kind, using Structure-from-Motion (SfM) to link detections across frames and correct the resulting double counts, and recent pipelines pair modern detectors with dedicated tracking modules for the same purpose \cite{rameshchandra2025}. In later work, Liu et al. \cite{liu2019} matched detections between frames with the Hungarian algorithm on top of a semantic SfM reconstruction, using a single camera. What these methods share is a reliance on densely sampled video: consecutive frames differ only slightly, so detections can be tracked through strong frame-to-frame correspondences, and the association is typically driven by 2D image-space cues such as bounding-box overlap or optical flow between adjacent views. Our setting is the opposite. To cover a tree from all sides within a limited view budget, the cameras are palced sparesly aroudn a full 360-degree range, so neighboring views often show markedly differet regions of the tree and share little image-space overlap. Frame-to-frame tracking inthe image plane is therefore not applicable, and the association must instead be established in 3D, through the reconstructed geometry rather than through 2D overlap. We adopt the Hungarian algorithm for this association as well, but apply it to correspondences derived from the 3D reconstructino rather than 2D bounding-box overlap between adjacent frames.

A third line of work removes double counting structurally by moving the count into three dimensions, so that each fruit exists once in a shared spatial representation. Häni et al. \cite{haeni2019} reconstructed rows of trees from both sides and merged the counts fromthe combined 3D models. FruitNeRF \cite{meyer2024fruitnerf} takes this idea furthest by integrating the counting directly into a neural scene representation, unifying reconstruction and segmentation that earlier methods had kept separate. Its drawback is that the neural representation must be optimized for each scene, which is computationally expensive; this cost is the immediate starting point of the present work. Newer methods reduce it by replacing the neural field with faster representations, suhc as the 3D Gaussian splatting used by FruitGaussian \cite{hong2025} and CountingFruit \cite{li2025}, but they remain tied to a per-scene optimization. A different way to avoid this optimization is to associate instance masks across views through a 3D representation rather than to optimize one; since this is the strategy of the present work, we treat it separately in the next section.

\section{Multi-View Instance Association and 3D Segmentation}
Associating 2D instance masks across many views into consistent 3D instances is a problem closely related to fruit counting, and one that general-purpose 3D scene-understanding methods address independently of any counting application. What distinguishes these methods is the criterion by which they decide that two masks seen in different views from the same object.

The most direct criterion is geometric overlap in an explicit reconstruction. Gené-Molá et al. \cite{genemola2020} segment fruit with Mask R-CNN, project the 2D masks into an SfM point could, associate detections by their overlap in 3D, and discard false detections with a support vector machine. The association is accurate but inherits the cost of the underlying SfM reconstruction, which the authors note is computationally intensive and unsuitable for real-time use.

A second criterion  replaces this geometric overlap with cross-view agreement. Maskclustering \cite{yan2024maskclustering} builds a global grpah whose nodes are the 2D masks and weights each edge by a \emph{view consensus rate}, the fraction of masks from other views that contain both endpoints; instances then emerge from iteratively clustering the masks that agree most strongly. The decision is thereby made globally, over all views at once, rather than only between neighboring frames, btu it rests on how masks co-occur across views rather than on explicit 3D correspondences.

The criterion our method builds on, introduced by CCGS \cite{hu2026ccgs}, derives the association directly from dense pointmaps. Pixel correspondences between views are found by minimizing the Euclidean distance between the predicted pointmaps; the overlap between two mass is redefined in therms of these correspondences; and the resulting matching cost ist solved with the Hungarian algorithm, with partial matched permitted. CCGS uses the associated masks to build a compact segmentation field for 3D Gaussian splatting, and it operates on image pairs from a DUSt3R reconstruction. We take this pointmap-based accociation and, instead, apply it to the joint multi-view reconstruction produced by VGGT, using it as the first stage of our counting pipeline.

\chapter{Fundamentals}
\section{VGGT}
The central component of the reconstruction stage in this work is VGGT (Visual Geometry Grounded Transformer), a feed-forward transformer introduced by Wang et al. \cite{wang2025vggt} and awarded the Best Paper Award at CVPR 2025. In this section, we first contextualize VGGT within the broader field of 3D reconstruction and then examine its input, data, internal pipeline, and outputs. Finally, we discuss the properties that make it a suitable foundation for our couting method.

\subsection{Context}
\label{sec:vggt_classification}
Reconstructing the three-dimensional structure of a scene from a set of two-dimensional images is a long-standing problem in computer vision \cite{hartley2004multiple}. The classical solution is a multi-stage pipeline consisting of Structure-from-Motion (SfM) and Multi-View Stereo (MVS).In the implementation by Schönberger and Frahm \cite{schoenberger2016sfm}, SfM estimates the camera poses as well as a sparsely populated point cloud by extracting and matching handcrafted features across the images, followed by an iterative optimization process known as bundle adjustment. A subsequent MVS step then densifies this reconstruction. While these methods are accurate and remain widely used, they rely on iterative optimization that is computationally expensive and can take several minutes to hours per scene to complete.

More recently, learning-based approaches have replaced parts of this pipeline with neural networks. Neural Radiance Fields (NeRFs) \cite{mildenhall2020nerf} represent a scene as a volumetric radiance function encoded in the weights of a neural network, while 3D Gaussian Splatting \cite{kerbl3Dgaussians} models it as a collection of colored Gaussian distributions. Both produce high-quality reconstructions but require separate optimization for each individual scene, which again incurs a substantial runtime cost per reconstruction.

A different line of work aims to eliminate per-scene optimization entirely by predicting geometry with a neural network. DUSt3R \cite{wang2024dust3r} pioneered this \emph{feed-forward} principle by regressing a dense pointmap from an image pair in a single network pass, followed by MASt3R \cite{mast3r_eccv24}, which added feature mapping. However, these methods operate on image pairs and still require a global alignment step to fuse the pairwise predictions into a consistent scene. VGGT \cite{wang2025vggt} belongs to this feed-forward family, but generalized it from pairs to arbitrary sets of images: It processes one, a few, or up to hundreds of views together in a single forward pass and predicts a globally consistent reconstruction without any subsequent optimization. This combination of speed and multi-view consistency motivates its use in this work, where fast reconstruction is a central goal.

\subsection{Input Data}
\label{sec:vggt_input}
The input for VGGT is a sequence of $N$ RGB images $(I_i)_{i=1}^{N}$ with $I_i \in \mathbb{R}^{3 \times H \times W}$, which capture the same static scene. Unluke the classical SfM pipeline, no camera parameters, calibration data, or prior pose information need to be provided: The model derives all geometric quantities directly from the image content. The number of views $N$ is flexible and can range from a single image to several hundred, with the only practical limitation being the available GPU memory, since all views are processed jointly \cite{wang2025vggt}.

\subsection{Pipeline}
\label{sec:vggt_pipeline}
Internally, VGGT maps the input sequence through tree conceptual stages to a set of dense 3D predictions: tokenization, alternating attention, and task-specific prediction heads. An overview is shown in Figure.
\missingfigure{vggt pipeline figure}

\paragraph{Tokenization.}
Each input image $I_i$ is first divided into non-overlapping image patches, which are encoded by a pre-trained DINOv2 backbone \cite{oquab2024dinov2} into a set of tokens. These tokens carry local appearance and structural information for their respective image regions. For each image, the patch totems are augmented with additional learnable tokens: a \emph{camera token}, which aggregates the information needed to predict the camera parameters of this image, as well as a set of \emph{register tokens}, which serve as auxiliary memory during the attention calculation an are discarded before the final prediction.

\paragraph{Alternating Attention.}
The core of the network is a transformer that processes all tokens of all images jointly. It consists of $24$ attention blocks that alternate between two types of self-attention. In \emph{frame-wise} self-attention, tokens attend only to other tokens of the same image, which preserves and refines the local structure of each image. In \emph{global} self-attention, all tokens attend to one another across all images, thereby exchanging information between views and establishing the correspondences that link the individual images into a single, coherent scene. By alternating between these two modes, the network baalnces local detail with global multi-view consistency. This alternating attention design is the central architectural mechanism of VGGT; the authors show that it is superior to both purely global attention and cross-attention alternatives \cite{wang2025vggt}.

\paragraph{Prediction Heads.}
After the attention blocks, the refined tokens are passed to specialized prediction heads. A dedicated \emph{camera head} processes the camera tokens to regress the intrinsic and extrinsic camera parameters of each view. A \emph{DPT head} \cite{Ranftl2021}, a Dense Prediction Transformer, upsamples the image tokens back to the input resolutions to generate the dense, pixel wise outputs, namely depth maps, point maps, tracking features, and the corresponding confidence maps. 

\subsection{Output}
\label{sec:vggt_output}
For each input iamge $I_i$, VGGT predicts the following quantities:
\begin{itemize}
    \item camera extrinsics and intrinsics of the view,
    \item a dense depth map together with a per-pixel confidence map,
    \item a dense point map together with a per-pixel confidence map, and
    \item dense feature maps used for tracking pixels across views.
\end{itemize}

The point map is the output most relevant to this work. It assigns a 3D scene coordinate $P_i(y) \in \mathbb{R}^{3}$ to every pixel $y$ of an image, thereby representing a dense, pixel-wise 3D reconstruction. Crucially, the point maps of all images are expressed in a single, common coordinate system: they are defined in the reference frame of the first camera, which serves as the world origin \cite{wang2025vggt}. Points observed in different views therefore already lie in a common 3D space without requiring an explicit alignment or registration step. This shared coordinate frame is the property our method builds upon, as it allows 2D instance masks from different views to be related to one another purely through their 3D coordinates. 

Each dense prediction is accompanied by a confidence map that estimates the pixel-wise reliability of the reconstruction. These confidence values provide a signal for identifying and filtering out unreliable regions, such as poorly observed or geometrically ambiguous areas \cite{wang2025vggt}.

\subsection{Advantages and limitations}
\label{sec:vggt_advantages}
The main advantage of VGGT for our application is its speed. Since the entire reconstruction is generated in a single forward pass, a scene can be reconstructed in a matter of seconds, whereas the SfM-based and scene-specific optimization approaches discussed in Section \ref{sec:vggt_classification} take minutes to hours. At the same time, according to the authors, VGGT matches or exceeds the accuracy of these optimization-based alternatives across multiple 3D tasks \cite{wang2025vggt}, and does so without any post-processing. Another practical advantage is that it operates directly on uncalibrated images and outputs camera parameters, depth, and point maps together, so that no separate calibration or psoe estimation step is required.

The fundamental limitationis that all input views are processed collectively via global self-attention, whose memory requirements increase with the number of tokens and, consequently, with the number of input views. The authors note that a naive implementation can become very memory intensive when dealing with a large number of tokens, so that in practice, the available GPU memory limits the number of images that can be reconstructed together \cite{wang2025vggt}. This constraint is directly relevant to our scenario, in which a fruit tree must be captured from all sides: The limited view budget forces a sparse distribution of cameras around the scene, which in turn weakens the correspondences between neighboring views. A number of follow-up works address this scaling limit, for example by processing long sequences in segments, as in VGGT-Long \cite{deng2025vggtlongchunkitloop} and VGGT-SLAM \cite{maggio2025vggtslam}, or by scaling up to thousands of images as in Fast3R \cite{Yang_2025_Fast3R}; Pi3 \cite{wang2025pi} additionally eliminates the dependence on a fixed reference view. These extensions are beyond the scope of this work, which uses the original VGGT model.

\section{Graph Theory}
Once VGGT has placed the pixels of all views into a common 3D coordinate frame, the reconstruction stage is follow by the task of deciding which 2D instance masks across the different views belong to the same physical apple. We formulate this assignment as a problem on a graph. In this section, we introduce the graph theory concepts used in our method and define them directly in the context of our problem. Building on this, we describe the "union-find" data structure, which we use to extract the corresponding instances from the graph.

\subsection{Basic concepts}
\label{sec:graph_basics}
A graph $G = (V, E)$ consists of a set of vertices (or nodes) $V$ and a set of edges $E \subseteq V \times V$, where each edge connects a pair of vertices. Two vertices are said to be \emph{adjacent} if they are connected by an edge. A \emph{path} is a sequence of vertices in which every consecutive pair of vertices is connected by an edge. A \emph{connected component} is a maximal subset of vertices in which every pair of vertices is connected by a path and which is not connected to any vertex outside the subset. Intuitively, a connected component is a group of nodes that are all directly or indirectly connected to one another, while being separated from the rest of the graph.

We apply these concepts to our association task as follows. Each \emph{node} corresponds to a single 2D instance mask in a single view, i.e. a pair $(i,m)$ consisting of an image index $i$ and an instance identifier $m$ within that image. An \emph{edge} between two nodes expresses the hypothesis that the two masks depict the same physical apple, observed from two different views. How these edges are established, namely by matching instances between image pairs, is described in Section \ref{sec:hungarian}. A \emph{connected component} then corresponds to a single apple: it groups together all instance masks across all views that were transitively linked to one another and therefore presumably depict the same fruit.

The last property is what makes the graph formulation useful for our purpose. Edges are only ever created between pairs of views, but an apple is typically visible in more than two images. By grouping nodes into connected components, an apple seen in, for example, five views is recovered as a single component of five nodes, even though no single matching step cmpared all five masks at once. THe transitivity of the relationship between connected components transforms the pairwise matches into a consistent mapping across multiple views. 

\missingfigure{simply illustration}

\subsection{Union-Find}
\label{sec:union_find}
Computing the connected components of a graph via a disjoint-set data structure is a standard technique \cite{cormen2009algorithms}. We use the \emph{Union-Find} data structure, also known as the disjoint-set data structure, which manages a partition of the nodes into disjoint sets and efficiently supports two operations: \texttt{find}, which returns a representative element for the set containing a given node, and \texttt{union}, which merges the two sets containing two given nodes into one.

Each set is stored as a tree, in which every node points to a parent node, and the root of the tree serves as the representative of the set. The \texttt{find} operation follows the parent pointers from a node up to the root. Two nodes belong to the same set if and only if the \texttt{find} operation returns the same root for both. The \texttt{union} operation joins two sets by attaching the root of one tree as a child of the root of the other.

In our method, we initialize the Union-Find structure with a singleton set per node, i.e., one set for each instance mask in each view. For each edge that is established during the matching step, we call \texttt{union} for its two nodes. After all edges have been processed, the sets of the Union-Find structure are exactly the connected components of the graph, and each set represents a counted apple instance.

A naive implementation of \texttt{find} can degenerate when the trees become tall, since each call then traverses a long chain of parent pointers. To avoid this, we apply \emph{path compression}: During a \texttt{find} operation, every node visited on the path to the root is immediately appended back to the root \cite{cormen2009algorithms}. This makes repeated lookups on the same set progressively cheaper and keeps the extraction of connected components efficient even for the large nzmber of instance nodes that arise across all views of a scene.

\section{Hungarian Matching}
\label{sec:hungarian}
The edges of the graph introduced in Section \ref{sec:graph_basics} are established by matching instance masks between pairs of veiws. We formulate this matching as a \emph{assignment problem} and solve it using the Hungarian algorithm \cite{kuhn1955hungarian, munkres1957algorithms}. Following the pointmap-based association of CCGS \cite{hu2026ccgs}, the cost of a match is derived from the overlap of two masks in 3D; the details of this cost are given in Chapter \ref{chap:method}, while this section introduces the assignment problem and the algorithm in general terms.

\subsection{The Assignment Problem}
The assignmetn problem is a classical combinatorial optimizataion problem. In its original formulation \cite{kuhn1955hungarian}, a set of $n$ workers is to be assigned $n$ jobs, given a numerical rating for the performance of each worker on each job, such that the total rating of the chosen assignment is optimal and no worker is assigned more than one job. Our matching task has the same structure: Instead of worker and tasks, we assign the instance masks of one view to those of another, and instead of a performance score, each candidate pair carries a cost value that reflects how likely it is that the two masks represent the same apple.

Consider two views $i$ and $j$, with instance masks indexed by $a \in \{1, \dots, n_i\}$ and $b \in \{1, \dots, n_j\}$ in views $j$. For every pair $(a, b)$, a cost $C_{ab} \geq 0$ expresses how unlikely it is that the two masks depict the same apple, where a lower cost indicates a stronger match. These costs form a cost matrix $C \in \mathbb{R}^{n_i \times n_j}$.

In its standard form, the linear sum assignment problem assumes an equal number of elements on both sides and assigns each element exactly once. Describing an assignment by binary variables $x_{ab} \in \{0, 1\}$, where $x_{ab} = 1$ means that mask $a$ is matched to mask $b$, it is written with equality constraints as $\sum_b x_{ab} = 1$ and $\sum_a x_{ab} = 1$, so that every element receives exactly one partner \cite{burkard2009assignment}. This assumption does not hold in our setting: The two views generally contain different numbers of masks, and a mask may have no counterpart at all when the apple it depicts is occluded in the other view. We therefore relax the equality constraints to inequalities, so that each mask is matched \emph{at most} once and unmatched masks are permitted. The resulting problem is 
\begin{align}
    \label{eq:assignment}
        \min_{x} \quad & \sum_{a} \sum_{b} C_{ab} \, x_{ab} \\
        \text{s.t.} \quad & \sum_{b} x_{ab} \leq 1 \quad \forall a, \nonumber\\
        & \sum_{a} x_{ab} \leq 1 \quad \forall b, \nonumber\\
        & x_{ab} \in \{0, 1\}. \nonumber
\end{align}
Because masks in view $i$ are only ever paired with masks in view $j$, and never with masks in the same view, the underlying structure is a \emph{bipartite} graph: its nodes split into the two disjoint sets of masks, and every candidate match connects one set to the other, with no edges inside either set. The cost matrix $C$ is the weighted adjacency matrix of this bipartite graph, and the assignment problem is the search for a minimum-cost bipartite matching.

The fact that $C$ is generally not square, and that some masks remain unmatched, is therefoer not a defect but the intended behavior: an apple that is visible in one view but occluded in the other should notbe assigned a partner. Each accepted match becomes an edge in the graph of Section \ref{sec:graph_basics}, linking the two corresponding nodes.

\subsection{The Hungarian Algorithm}
The difficulty of the assignment problem is that the number of possible assignments grows factorially with the number of masks: for two views with $n$ masks each, there are $n!$ distinct one-to-one assignments \cite{munkres1957algorithms}, so enumerating them and picking the cheapest is infeasible for all but the smallest cases. A tempting shortcut is a greedy strategy in which each mask is assigned to its individually best partner. Such a choise is locally optimal but can be globally suboptimal: two masks in view $i$ might both prefer the same mask in view $j$, which violates the one-to-one requirement, and resolving the conflict greedily can force a worse overall assignment than necessary.

The Hungarian algorithm avoids both problems. Introduced by Kuhn \cite{kuhn1955hungarian}, it is named for the the underlying results of the Hungarian mathematicians D. König and E. Egerv\'ary on which it is based, and it computes a globally optimal assignment without enumerating all possibilities. Munkres \cite{munkres1957algorithms} later gave a variant of the method together with a proof that it terminates after a number of operations bounded by a polynomial in $n$, rather than the factorial cost of exhaustive search; because of this joint contribution the method is also konwn as the Kuhn-Munkres algorithm. It therefore solves the problem in \eqref{eq:assignment} exactly and efficiently. In our implementation, we obtain the optimal assignment using a standard solver for the linear assignment problem.

\subsection{Partial Matching}
The optimal assignment minimizes the total cost, but it does not by itself guarantee that the matched pairs are meaningful: the algorithm will continue to pair two masks whose overlap is negligible if doing so reduces the total cost. To prevent such spurious matches from becoming edges, we allow an assignment to be an edge only if its overlap exceeds a threshold and discard the remaining assignments. This corresponds to allowing \emph{partial matches}, in which not every mask of the smaller view has to be matched, and follows the association strategy of CCGs\cite{hu2026ccgs}. The precise overlap measure and threshold are described in Chapter \ref{chap:method}.


\section{RANSAC}
\section{Clustering}

% German: Methode
\chapter{Method}
\label{chap:method}

% German: Ergebnisse
\chapter{Results}
Describe your qualitative and quantitative evaluations and present your results.
Discuss the limitations of your work.

% German: Zusammenfassung
\chapter{Conclusions}
Give a brief summary of your work and your contributions.
Afterwards, briefly outline potential opportunities for future work.

\listoftodos
\listoffigures
\listoftables
\lstlistoflistings

% bibliography and declaration are required
\printbibliography
\declaration

\end{document}
